\originalchapter{半导体、半经典运动与量子Hall效应}
\label{ch:transport}
\source{8.511第11--14讲；8.513第1--6讲的先修连接。}
\originalsection{有效质量、掺杂与激子}
在能带极值附近展开$E_n(\vct k_0+\vct q)=E_0+\hbar^2q_i(m^{*-1})_{ij}q_j/2$，故$(m^{*-1})_{ij}=\hbar^{-2}\partial_i\partial_j E$。它是张量，可有不同主轴。价带顶的电子曲率为负；描述缺少的电子时取带正电$+e$的空穴，可获得正的空穴有效质量。

浅施主的长距离吸引势为$-e^2/(4\pi\epsilon_0\epsilon_r r)$。各向同性有效质量近似把氢原子的$m$换成$m^*$、真空介电常数换成$\epsilon_0\epsilon_r$，于是$a^*=\epsilon_r(m/m^*)a_B$、束缚能$E_b=(m^*/m)13.6\,\mathrm{eV}/\epsilon_r^2$。当$a^*$远大于晶格常数时这个包络近似自洽；深杂质的中央胞修正不属于它。电子与空穴形成Wannier激子时把$m^*$换成约化质量$\mu^{-1}=m_e^{*-1}+m_h^{*-1}$，同样要求弱束缚和大半径。间接带隙半导体的光学跃迁还需声子提供晶体动量，这解释电学带隙与光学吸收边之间的区别。
\originalsection{波包方程与轨道量子化}
先忽略Berry曲率，窄波包的中心满足$\dot{\vct r}=\hbar^{-1}\nabla_kE_n$、$\hbar\dot{\vct k}=-e(\vct E+\dot{\vct r}\times\vct B)$。条件是外场在原胞尺度上缓慢变化，且带间跃迁小。均匀电场使$k(t)=k(0)-eEt/\hbar$，一维紧束缚速度随$k$周期变化，所以位移有Bloch振荡，频率$\omega_B=eEa/\hbar$。实际散射使振荡可见性要求$\omega_B\tau\gtrsim1$，强场又可能诱发Zener跨带跃迁。

在$\vct E=0$且$\vct B=B\hat z$时，$k_z$与能量守恒，电子沿$k$空间等能截线运动。对电荷$-e$的电子，$\hbar\dot{\vct k}=-e(\dot{\vct r}\times B\hat z)$；两边左叉乘$\hat z$后消去速度，得$\dd\vct r_\perp=-(\hbar/eB)\hat z\times\dd\vct k_\perp$。这确定了轨道的取向，实空间面积与动量空间面积之比为$(\hbar/eB)^2$。闭合轨道的作用量量子化给Onsager规则
\begin{equation}
A_k(E,k_z)=\frac{2\pi eB}{\hbar}(n+\gamma),\qquad m_c=\frac{\hbar^2}{2\pi}\frac{\partial A_k}{\partial E}.
\end{equation}
无Berry相位的简单抛物线$\gamma=1/2$。$m_c$控制回旋频率$eB/m_c$，不等于任意主轴质量。电子能级经过$E_F$时态密度振荡，$1/B$周期为$\Delta(1/B)=2\pi e/(\hbar A_{\rm ext})$，其中$A_{\rm ext}$是垂直磁场的极值截面积。改变磁场方向便能探测费米面形状。热展宽和有限寿命削弱振荡；磁击穿是强场下相邻轨道之间跨越小能隙的隧穿。
\originalsection{Landau态与边缘态}
二维抛物线电子取$\vct A=-By\hat x$，$p_x$守恒。哈密顿量在$y_0=\hbar k_x/(eB)$附近等价于频率$\omega_c=eB/m$的谐振子，因此$E_n=\hbar\omega_c(n+1/2)$。把导引中心限制在宽$L_y$内，$k_x$间距$2\pi/L_x$给每Landau级每自旋简并$N_\phi=eBL_xL_y/h$。磁长$\ell_B=\sqrt{\hbar/(eB)}$决定波函数宽度，且$N_\phi=A/(2\pi\ell_B^2)$。

若$n$个自旋分辨Landau级填满，体态有能隙。在边界，限制势$V(y)$抬升能级，$E_n(k_x)\simeq\hbar\omega_c(n+1/2)+V(y_0)$，边缘群速度$\hbar^{-1}\dd E/\dd k_x=V'(y_0)/(eB)$。两侧边缘方向相反；一条边同向传播的通道缺少同位置反向态，弱无序不能直接造成后向散射。

把二维电子放在圆柱或环上，沿孔缓慢插入磁通$\Phi_0=h/e$。规范变换把插通量等价成边界相位$2\pi$变化，每条填满带泵送一个电子从一边到另一边。若有$n$条通道，$\Delta Q=ne$；由Faraday电动势积分$\int\mathcal E\dd t=\Phi_0$得$\sigma_{xy}=\Delta Q/\Phi_0=ne^2/h$，符号由电荷与磁场方向约定。此论证要求绝热过程中体能隙或迁移率隙不闭合。实际平台的宽度需要局域态容纳额外粒子，纯净固定密度Landau级计数本身不能给整个平台。
\originalsection{Berry修正的连接}
Bloch本征矢随$\vct k$变化时，定义$\mathcal A_n(\vct k)=\ii\avg{u_{n\vct k}|\nabla_k u_{n\vct k}}$及$\vct\Omega_n=\nabla_k\times\mathcal A_n$。波包作用量含$\hbar\mathcal A_n\cdot\dot{\vct k}$，Euler--Lagrange方程使速度多出$-\dot{\vct k}\times\vct\Omega_n$。该项是异常速度，来源于态空间几何而非额外力。第三编会从绝热演化推导Berry相位与Chern数，再说明没有外磁场的Chern绝缘体如何也有量子化Hall响应。
